\documentclass[11pt]{article}
\usepackage[margin=1in]{geometry}
\usepackage{amsmath,amssymb,amsthm}
\usepackage{booktabs}
\usepackage{natbib}
\usepackage[hidelinks]{hyperref}
% Shared notation for the differential-geometric SN Ia light-curve project.
% Input this file from every scientific document so notation stays consistent.
%
% Requires: amsmath, amssymb

% --- Number sets -------------------------------------------------------
\newcommand{\Real}{\mathbb{R}}
\newcommand{\Rn}{\mathbb{R}^{n}}
\newcommand{\Rtwo}{\mathbb{R}^{2}}
\newcommand{\Rthree}{\mathbb{R}^{3}}

% --- Differential geometry ---------------------------------------------
\newcommand{\dd}{\mathrm{d}}
\newcommand{\curve}{\gamma}                 % the light-curve curve
\newcommand{\arc}{s}                        % arclength
\newcommand{\curv}{\kappa}                  % curvature
\newcommand{\tors}{\tau}                    % torsion (n >= 3 only)
\newcommand{\frenet}[1]{\mathbf{E}_{#1}}    % Frenet frame vector
\newcommand{\tang}{\mathbf{T}}              % unit tangent
\newcommand{\norml}{\mathbf{N}}             % unit normal
\newcommand{\arclen}{L}                     % total arclength of the window

% --- Model parameters ---------------------------------------------------
\newcommand{\shape}{\theta}                 % shape latents
\newcommand{\shapeone}{\theta_{1}}
\newcommand{\shapetwo}{\theta_{2}}
\newcommand{\tscale}{\sigma}               % timing scale
\newcommand{\tmax}{t_{\mathrm{max}}}        % time origin
\newcommand{\warpone}{a_{1}}                % quadratic timing corrections
\newcommand{\warptwo}{a_{2}}
\newcommand{\dm}{\mu}                       % normalization / distance modulus
\newcommand{\colr}{c}                       % colour translation (variant only)

% --- Photometry ---------------------------------------------------------
\newcommand{\magg}{m_{g}}
\newcommand{\magr}{m_{r}}
\newcommand{\fluxg}{f_{g}}
\newcommand{\fluxr}{f_{r}}
\newcommand{\zp}{\mathrm{zp}}
\newcommand{\phase}{p}                      % rest-frame phase
\newcommand{\zred}{z}

% --- Objects and abbreviations -----------------------------------------
\newcommand{\snia}{SN\,Ia}
\newcommand{\sneia}{SNe\,Ia}
\newcommand{\ztf}{ZTF}
\newcommand{\salt}{SALT2}
\newcommand{\NN}{\mathrm{NN}}

% --- Operators ----------------------------------------------------------
\DeclareMathOperator{\SO}{SO}
\DeclareMathOperator*{\argmin}{arg\,min}
\newcommand{\deriv}[2]{\frac{\dd #1}{\dd #2}}
\newcommand{\pderiv}[2]{\frac{\partial #1}{\partial #2}}
\newcommand{\abs}[1]{\left| #1 \right|}
\newcommand{\norm}[1]{\left\| #1 \right\|}


\title{A Differential-Geometric Model for \\ Type Ia Supernova Multi-Band Light Curves}
\date{\today}

\begin{document}
\maketitle

\section{Motivation}

The standard empirical treatment of \snia\ light curves, exemplified by \salt,
models the spectral energy distribution as a linear expansion about a mean
sequence,
\begin{equation}
  F(p,\lambda) = x_{0}\left[M_{0}(p,\lambda) + x_{1}M_{1}(p,\lambda) + \dots\right]
                 \exp\!\left[\colr\,\mathrm{CL}(\lambda)\right],
\end{equation}
with per-supernova parameters $(x_{0}, x_{1}, \colr, t_{0})$: a description of a
\emph{surface} over phase and wavelength, from which the light curve in any band
follows by integration against a filter response. The parameter $x_{1}$ is read as a
``stretch'', but it does not act as a reparameterisation of time --- it sets the
amplitude of a fixed additive component --- so shape and timing are entangled in one
number.

The geometric alternative starts elsewhere. A supernova observed in $n$ bands
supplies a point in $\Rn$ at each epoch, and over the event those points trace a
curve. The curve is independent of when it was sampled, and by the fundamental
theorem of curves it is determined, up to a rigid motion, by $n-1$ scalar invariants
of arclength. Shape then lives entirely in those invariants and timing entirely in
the arclength map, since a reparameterisation of time changes how fast a supernova
moves along its curve but does not change the curve.

That separation is exact as a statement about curves and is not thereby a guarantee
about fits; Section~\ref{sec:gauge} gives the degeneracy that limits it in practice.

\section{The model}

The objective of the model is to represent rest-frame multi-band light curve template as a function of phase
\begin{equation}
  \curve(\phase; \theta) = \bigl(M_{1}(\phase; \theta), \dots, M_{n}(\phase; \theta)\bigr) \in \Rn ,
  \label{eq:curvedef}
\end{equation}
where $\theta$ are per-supernova parameters that capture the diversity of SNe~Ia.

The model represents light curves using three objects: curves in magnitude space, the positions and orientations of those curves on the magnitude coordinate system, and the traversal of supernovae along those curves.  Within the framework of differential geometry,
curves are specified by curvature, positions and orientations by boundary conditions, and traversal with 
unit arclength as a function of phase.  These concepts will be described in Sections ~\ref{sec:curve}, \ref{sec:place}, and
Section~\ref{sec:traversal} respectively. 

The primary configuration is \ztf\ $g$ and $r$ \citep{Bellm2019}, so $n=2$ and
$\curve = (M_g, M_r)$; a subset of the low-redshift sample also carries usable
$i$-band photometry, giving $n=3$ on fewer objects (Section~\ref{sec:samples}). The
implementation is written for general $n$ for that reason, not as speculative
generality.

We work in magnitude space, with the Euclidean metric
$\mathrm{diag}(1,1)$. Both are choices, argued in Section~\ref{sec:magspace}: the
space because distance and dust (to first order) act there as translations, to which curvature
is blind; the metric because none is physically preferred, so one must be declared.
The likelihood, by contrast, is evaluated in native flux
(Section~\ref{sec:likelihood}).

\subsection{The curve}
\label{sec:curve}

Reparameterise by arclength $\arc$, defined by
$\dd \arc = \norm{\dd\curve/\dd \phase}\,\dd \phase$, so that
$\norm{\curve'(\arc)} = 1$. A unit-speed curve in $\Rn$ admits a Frenet frame
$\frenet{1},\dots,\frenet{n}$ obeying
\begin{align}
  \frenet{1}' &= \curv_{1}\frenet{2}, \nonumber\\
  \frenet{i}' &= -\curv_{i-1}\frenet{i-1} + \curv_{i}\frenet{i+1},
                 \qquad 1 < i < n, \label{eq:frenet}\\
  \frenet{n}' &= -\curv_{n-1}\frenet{n-1}, \nonumber
\end{align}
together with $\curve' = \frenet{1}$. The $n-1$ functions
$\curv_{1}(\arc),\dots,\curv_{n-1}(\arc)$ determine the curve up to a rigid motion of
$\Rn$. For $n=2$ this collapses to a single invariant, the signed planar curvature.
The indexed notation is conventionally dropped in the plane, and we drop it
throughout: write $\tang \equiv \frenet{1}$ for the unit \emph{tangent},
$\norml \equiv \frenet{2}$ for the principal \emph{normal}, and
$\curv \equiv \curv_{1}$ for the \emph{curvature}. In this notation
\eqref{eq:frenet} reads
\begin{equation}
  \curve' = \tang, \qquad
  \tang' = \curv\,\norml, \qquad
  \norml' = -\curv\,\tang .
  \label{eq:frenet2d}
\end{equation}
In the plane $\norml$ is fixed by rotating $\tang$ through $+90^{\circ}$. This agrees
with the classical principal normal $\curve''/\norm{\curve''}$ wherever
$\curv \neq 0$ and, unlike it, remains defined where $\curv$ vanishes --- a
distinction that becomes load-bearing below.
For $n=3$ a second invariant appears, the torsion $\tors = \curv_{2}$, measuring
departure from the instantaneous osculating plane.

Because the curvature functions $\curv_{i}(\arc)$ determine the curve up to a rigid
motion, modelling the shape of a supernova's track in the magnitude plane means
modelling its $\curv_{i}(\arc)$. A neural network supplies them, and a per-supernova
latent vector $\shape \in \Real^{K}$ lets them vary from one supernova to the next,
\begin{equation}
  \curv_{i}(\arc; \shape) = \NN_{i}\bigl(\arc, \shape\bigr),
  \qquad i = 1,\dots,n-1 .
  \label{eq:kappanet}
\end{equation}

The domain is a bounded interval of arclength,
\begin{equation}
  \arc \in [\arc_{\min}, \arc_{\max}],
  \label{eq:domain}
\end{equation}
held in configuration and taken wide enough that the traversal of
Section~\ref{sec:traversal} stays inside it across the selection window for every sampled
parameter value --- a condition to be checked rather than assumed. Nothing is imposed on
$\curv_{i}$ within the interval: the invariants are learned wherever the fit can see them.

Vanishing curvature is nevertheless benign for $n = 2$, by a property of
\eqref{eq:frenet2d} rather than by a modelling choice. The planar Frenet system is a
rotation ODE, regular at $\curv = 0$, because $\norml$ is \emph{defined} as the
$90^{\circ}$ rotation of $\tang$ and not as $\curve''/\norm{\curve''}$, which is undefined
there. Integrating the system rather than differentiating the curve is therefore the
definition that survives an inflection --- which the $(\magg, \magr)$ track may approach
incidentally near the secondary maximum. For $n \ge 3$ the protection does not extend to
the second invariant: where $\curv_{1} = 0$ there is no osculating plane, so $\tors$ is
not merely ill-conditioned but undefined, the frame must be carried by parallel transport
rather than by the classical construction, and $\tors$ must be reported as unidentifiable
wherever $\curv_{1}$ approaches zero.

\subsection{Where the curve sits}
\label{sec:place}

Placement is the initial data of the Frenet system and nothing more. In the plane
\eqref{eq:frenet2d} closes on $(\curve, \tang)$ --- $\norml$ is the fixed rotation of
$\tang$ and is never independent data --- so the system is a first-order linear ODE whose
solution is determined by $\curve$ at one arclength together with the direction of $\tang$
at one arclength. That is three numbers, which is exactly $\dim\SE(2)$, the group the
fundamental theorem of curves quotients by. Placement therefore needs no boundary
condition at infinity and no appeal to the pre-explosion state; it is three constants, and
the only question is which of them are free per supernova.

For $n \ge 3$ the count is $n + n(n-1)/2 = \dim\SE(n)$ and a tangent direction is not
sufficient data, since the system evolves the whole frame: the initial data are $\curve$
together with an element of $\SO(n)$, six numbers for $n = 3$.

\paragraph{The translation is two per-supernova parameters.} The geometry treats its two
directions quite differently, though $\curv$ is blind to both by \eqref{eq:transinv}.
Along $(1,1)$ lies distance: luminosity distance, absolute luminosity and peculiar
velocity each shift every band by the \emph{same} amount, so this direction is fixed by
physics rather than chosen. The second direction is reddening, and which direction that is
--- the reddening vector, set by the extinction law and the mean SED --- no per-supernova
physics fixes. Rather than assume it, the model carries it as a single sample-wide unit
vector $\hat{\mathbf{e}}_{c}$, the two-band analogue of \salt\ fitting $\beta$ rather than
adopting a value of $R_{V}$, whose fixed-$R_{V}$ idealisation
$\hat{\mathbf{e}}_{c} = (1,-1)/\sqrt{2}$ appears in \eqref{eq:transcol}. The per-supernova
offsets are then $\dm$ along $(1,1)$ and $\colr$ along $\hat{\mathbf{e}}_{c}$; the two
directions are independent, so they span the translation freedom in full and no third
per-supernova placement parameter exists to be added.

That completeness settles the status of $\hat{\mathbf{e}}_{c}$, which is easily over-read
as fitted at every rung. Wherever $\colr$ multiplies a constant vector, $\dm$ and $\colr$
already span $\Rtwo$ for \emph{any} $\hat{\mathbf{e}}_{c}$ not parallel to $(1,1)$, so
changing the direction is undone exactly by an invertible relabelling of each supernova's
$(\dm, \colr)$: the likelihood is invariant, and the direction is an exactly flat
direction of it. Below rung L2c $\hat{\mathbf{e}}_{c}$ is therefore a gauge choice and not
a measurement. It acquires content only at L2c, where $\colr$ also multiplies the
phase-varying $\delta\dustdir(\arc)$ of \eqref{eq:dustsplit}, which no translation
absorbs and which ties the amplitude to a definite direction.

\paragraph{The rotation is fixed by convention.} The invariants and the orientation are
not independent: $\curv(\cdot\,;\shape)$ paired with an initial tangent angle $\psi$
generates the same curve as $\curv(\cdot + a;\shape)$ paired with
$\psi + \int_{0}^{a}\curv$, so the network's arclength origin and the frame carry one
joint redundancy, and one condition removes it. The condition adopted places the arclength
origin at the principal $g$-maximum,
\begin{equation}
  \deriv{\magg}{\arc}\bigg|_{\arc = 0} = 0
  \qquad\text{with}\qquad
  \tang(0) = (0,-1), \quad \curv(0) > 0 ,
  \label{eq:frameanchor}
\end{equation}
the sign making $\arc = 0$ the brightest point in $g$ --- a \emph{minimum} of $\magg$ ---
rather than the faintest, with $\tang(0) = (0,+1)$, $\curv(0) < 0$ the reflected
alternative. This costs no generality: every planar curve that turns has a point of
vertical tangent, and calling it $\arc = 0$ relabels rather than restricts. Nor does it
constrain which band peaks first, since a stationary point of $\magr$ is reachable on
either side of the origin according to how far the tangent has turned.

The bookkeeping then closes differently in the two ambient spaces, and this is what the
earlier open question about the $\Rthree$ frame amounts to. For $n = 2$ the single
rotational number and the single redundancy cancel: orientation is \emph{gauge}, and no
rigid degree of freedom survives to be fitted. For $n \ge 3$ they do not cancel ---
$n(n-1)/2$ rotational numbers against one redundancy --- so $n(n-1)/2 - 1$ orientation
parameters survive, two for $n = 3$, shared across the sample and fitted globally rather
than per supernova. Not a missing condition, then, but a definite and small number of
sample-wide parameters.

Two conveniences follow immediately for $n = 2$. The traversal of
Section~\ref{sec:traversal} has only to map $\tmax$ onto $\arc = 0$, with no
$\shape$-dependent root to locate; and normalising the template to $\curve(0) = 0$ makes
the per-supernova offset $\dm$ of \eqref{eq:forward} along $(1,1)$ \emph{exactly} the peak
$g$ magnitude, while $\colr$ carries the colour at that epoch. Both are conventions, and
with \eqref{eq:frameanchor} they exhaust the three constants counted above.

\paragraph{Zero flux is an output, not an input.} In every band $f \to 0$ at both
extremes, hence $m \to +\infty$: both ends of the track escape to infinity, so the curve
is a hairpin rather than an arc, and the tangent reverses between the ends, turning
through approximately $\pi$. None of this is imposed. Both asymptotes lie outside
\eqref{eq:domain} and neither is observed, so $\curv \to 0$ at the extremes, the total
turning $\int\curv\,\dd\arc \approx \pi$, and the offset between the asymptotic rays as
the terminal colour are all \emph{consequences} of a fitted $\curv$ --- qualitative
validation numbers on its extrapolation, never constraints. Both translations leave
$\curv$ untouched by \eqref{eq:transinv}, with one consequence for colour worth stating
precisely: with $\colr$ free, neither the early nor the terminal colour is predicted in
absolute terms, only their \emph{difference}, since $\colr$ displaces the whole curve and
cancels from it. The model predicts colour \emph{evolution} and leaves the colour zero
point free --- structurally the division \salt\ makes between a fixed colour law and a
fitted $\colr$. One regularity condition on the placement does remain.

The band maxima must not coincide. If all bands are stationary at the same epoch then
$\norm{\dd\curve/\dd\phase} = 0$, hence $\dd\arc/\dd t = 0$: arclength ceases to
increase, so no monotone time-to-arclength map exists and the unit-speed
parameterisation fails, the curve having a stationary rather than a regular point. Non-coincident band maxima is thus a
regularity condition on the data, distinct from the vanishing-$\curv$ degeneracy of
Section~\ref{sec:curve}, which concerns turning, not speed, and leaves an inflection
perfectly regular.

\subsection{The traversal}
\label{sec:traversal}

Photometry is recorded in observer-frame MJD while the curve is a rest-frame object,
so observer-frame intervals are divided by $1+\zred$ before anything else happens.
Since $\zred$ is measured for every object this is a \emph{known} kinematic rescaling,
divided out deterministically rather than fitted --- the same logic applied to
\texttt{mwebv} in Section~\ref{sec:samples}: known effects are removed, not absorbed
into free parameters. It keeps $\tscale$ a genuine rest-frame stretch rather than a
repackaging of time dilation.

Timing is then a monotone map from observer time to arclength. With $u$ the phase in units
of the timing scale $\tscale$, itself in rest-frame days,
\begin{equation}
  \arc(t) = u\left[1 + \warpone u + \warptwo u^{2}\right],
  \qquad
  u = \frac{t - \tmax}{\tscale\,(1+\zred)},
  \label{eq:reparam}
\end{equation}
with $\deriv{\arc}{t} > 0$ required across the window. At $t = \tmax$ one has $u = 0$ and
$\arc = 0$, which by \eqref{eq:frameanchor} is the template's $g$-maximum; and since
$\dd\arc/\dd t > 0$ implies $\dd\magg/\dd t = 0 \iff \dd\magg/\dd\arc = 0$, the fitted
$\tmax$ \emph{is} the observer-frame epoch of $g$-band maximum. Nothing is added to
enforce that reading --- the gauge of Section~\ref{sec:place} enforces it, and it is the
reason that gauge was chosen. As before, $\tmax$ is free for every supernova and is never
taken from the archive's \salt\ $t_{0}$.

Setting $\warpone = \warptwo = 0$ recovers the affine map, the direct analogue of
classical stretch \citep{Phillips1993}; the quadratic and cubic terms permit asymmetric
warping between rise and decline. Positivity is definitional rather than regulatory ---
$\arc$ is arclength, so $\dd\arc/\dd t = \norm{\dot{\curve}}$ is strictly positive
wherever the supernova is moving --- but it is not automatic in $(\warpone, \warptwo)$ and
must be checked over the sampled values, together with the requirement that $\arc(t)$ stay
inside \eqref{eq:domain}.

\paragraph{What the traversal deliberately does not do.} Nothing above turns the supernova
on. Predicted flux is positive at every arclength of \eqref{eq:domain}, so epochs before
explosion are not described by the model: they are excluded from the likelihood rather
than fitted, and no explosion epoch appears among the parameters. This is a restriction of
scope, not an oversight, and the reason it is a large one is worth recording. A turn-on
needs $\arc \to -\infty$, because the curve is unit-speed and
$\norm{\curve(\arc) - \curve(\arc_{0})} \le \abs{\arc - \arc_{0}}$ makes $f = 0$
unreachable at finite arclength; that in turn needs a semi-infinite domain, a traversal
singular at the explosion epoch, and curvature prescribed over a region no photometry
constrains. Those pieces fit together, and their payoff is a predicted early power-law
rise, but they are a candidate later rung rather than part of the model defined here, and
the deliverable does not rest on them.

Phase zero is conventionally set at
rest-frame $B$ maximum, but locating $B$ maximum requires K-correcting into a filter
that was not observed, hence an SED, which a two-band model does not have. The anchor
here is the maximum of observed \ztf\ $g$. Since $g$ and $B$ are different filters,
$\tmax$ and a $B$-maximum epoch need not coincide, and the offset is not predictable
from effective wavelength: peak epoch is not monotonic in wavelength from UV to NIR.
When \salt\ is used as an external benchmark the distribution of $\tmax - t_{0}$ must
be \emph{measured} --- neither its centre nor its sign assumed --- and its scatter, not
its mean, is the check on timing recovery.

\subsection{From curve to observables}
\label{sec:likelihood}

Given latents, the predicted magnitude in band $X$ at observer-frame epoch $t$ is,
with $\arc(t)$ from \eqref{eq:reparam} and $\curve$ from \eqref{eq:frenet},
\begin{equation}
  m_{X}^{\mathrm{pred}}(t)
    = \curve_{X}\bigl(\arc(t)\bigr) + \dm
      + \colr\,\dustdir_{X}\bigl(\arc(t)\bigr) ,
  \label{eq:forward}
\end{equation}
and where host reddening displaces the curve by
\begin{equation}
  \curve(\arc) \longmapsto \curve(\arc) + \colr\,\dustdir(\arc),
  \qquad
  \dustdir(\arc) \ \text{a fixed vector-valued function of arclength, not a constant,}
  \label{eq:dustcurve}
\end{equation}
with the \emph{phase variation} of $\dustdir$ precomputed once from an extinction law
\citep{Fitzpatrick1999} and a template SED \citep{Hsiao2007}, its mean direction the
fitted global $\hat{\mathbf{e}}_{c}$ of Section~\ref{sec:place}, and its amplitude
$\colr$ free per supernova. That $\dustdir$ \emph{must} depend on $\arc$, and what
follows from it, is Section~\ref{sec:phasedust}.

The amplitude $\colr$ is present at every rung; what changes along the ladder is the
function it multiplies. Splitting $\dustdir(\arc) = \bar{\dustdir} + \delta\dustdir(\arc)$
with $\bar{\dustdir} = \langle \dustdir \rangle$ and $\langle \delta\dustdir \rangle = 0$,
the baseline takes $\dustdir \to \bar{\dustdir}$, a constant vector along the fitted
global reddening direction $\hat{\mathbf{e}}_{c}$ (Section~\ref{sec:place}): a pure
translation, invisible to $\curv$, and exactly \salt's colour parameter. Rung L2c
restores the full $\arc$-dependence, importing the phase variation
$\delta\dustdir(\arc)$ from the template while $\hat{\mathbf{e}}_{c}$ continues to set
the mean direction. The two differ only by $\colr\,\delta\dustdir(\arc)$,
which is second order --- small in $\colr$ and again in $\abs{\delta\dustdir}/\abs{\bar{\dustdir}}$
--- and being phase-dependent is the part that cannot be folded into $\dm$ or
$\bar{\dustdir}$ after the fact, so it is the part that contaminates $\shape$. L2c
therefore adds no parameter; it tests whether that second-order term matters.

Equation~\eqref{eq:forward} is converted to flux on the instrumental system,
\begin{equation}
  f_{X}^{\mathrm{pred}}(t)
    = 10^{-0.4\left(m_{X}^{\mathrm{pred}}(t) - \zp_{X}(t)\right)} .
  \label{eq:tofluxdef}
\end{equation}
The magnitudes are finite throughout \eqref{eq:domain}, so this map is an ordinary change
of variable and carries no boundary case; it is here because the errors are Gaussian on
its right-hand side, not on its left (Section~\ref{sec:magspace}).

Errors are Gaussian in flux, so the comparison is made there.
For band $X$ let $\mathbf{r}_{X}$ stack the residuals
$f^{\mathrm{obs}}_{X,k} - b_{X} - f_{X}^{\mathrm{pred}}(t_{k})$ over epochs $k$,
where $b_{X}$ is the reported difference-imaging zero level. Then
\begin{equation}
  -2\ln\mathcal{L} = \sum_{X}
    \mathbf{r}_{X}^{\!\top} C_{X}^{-1} \mathbf{r}_{X},
  \qquad
  C_{X} = \mathrm{diag}\bigl(\sigma_{X,k}^{2}\bigr)
          + \bigl(\sigma^{0}_{X}\bigr)^{2}\,\mathbf{1}\mathbf{1}^{\!\top},
  \label{eq:chi2}
\end{equation}
with $\sigma_{X,k}$ the reported flux error of epoch $k$ and $\sigma^{0}_{X}$ the
uncertainty on $b_{X}$. Here $\sigma$ denotes measurement error throughout and never the
timing scale, which is $\tscale$; and $\varsigma$ is reserved for the classical stretch
factor of Section~\ref{sec:salt}.

The rank-one term is not decoration. The zero level $b_{X}$ is estimated once per light
curve, so its error is \emph{common} to every epoch of that band: negligible beside a
bright detection, but comparable to the signal across the faint epochs late in the window,
where a diagonal $C_{X}$ would make their joint constraint appear to tighten as $N^{-1/2}$
though the true floor is $\sigma^{0}_{X}$ and does not shrink. Since $b_{X}$ and
$\sigma^{0}_{X}$ are both reported, the correct treatment costs no parameter: subtract
$b_{X}$ and invert $C_{X}$ in closed form by Sherman--Morrison. The same rule as for
$1+\zred$ and \texttt{mwebv}.

Two windows must be kept apart, having previously shared one symbol. \emph{Model support}
is the arclength interval \eqref{eq:domain}, which the traversal must not leave. The
\emph{selection window} is $\phase \in [-15,+40]$~d of Section~\ref{sec:samples}, and it
also delimits the likelihood epoch set: epochs before explosion are not described by the
model (Section~\ref{sec:traversal}) and are excluded rather than fitted.

Since $\tmax$ is free, the selection window is a function of the fitted parameters: it
spans $57.75$~d of observer time at $\zred = 0.05$ against $55$~d at $\zred = 0$, and
slides as $\tmax$ moves.

\subsection{Per-supernova parameters}

Eight per supernova, each belonging to one layer above:
\begin{center}
\begin{tabular}{llll}
\toprule
Layer & Symbol & Role & Geometric character \\
\midrule
Curve (\ref{sec:curve})       & $\shapeone, \shapetwo$ & shape & condition $\curv$ \\
Placement (\ref{sec:place})   & $\colr$                & colour & translation $\parallel \hat{\mathbf{e}}_{c}$ \\
Traversal (\ref{sec:traversal}) & $\tscale$            & timing scale (rest-frame d) & reparameterisation \\
                              & $\warpone, \warptwo$   & timing warp & reparameterisation \\
Observables (\ref{sec:likelihood}) & $\tmax$           & epoch of $g$-maximum (observer MJD) & reparameterisation \\
                              & $\dm$                  & peak $g$ magnitude & translation $\parallel (1,1)$ \\
\bottomrule
\end{tabular}
\end{center}
Placement costs exactly one per-supernova parameter, and only in position: the orientation
is gauge for $n = 2$ by \eqref{eq:frameanchor}, and the reddening direction
$\hat{\mathbf{e}}_{c}$ is one sample-wide quantity, not a per-supernova cost. The component
of position along $(1,1)$ is normalised away on the template, leaving $\dm$ to carry it
per supernova, while the component along $\hat{\mathbf{e}}_{c}$ carries $\colr$. $\tmax$
enters the traversal
mechanically through \eqref{eq:reparam} but is interpreted as an observable. In the
baseline $\colr$ multiplies the constant $\bar{\dustdir}$ and is a genuine translation;
at rung
L2c it multiplies the full $\dustdir(\arc)$ and is then not a motion of the ambient
space at all, moving different points of the curve by different amounts
(Section~\ref{sec:phasedust}). The parameter count is the same either way.

For comparison \salt\ uses four, $(x_{0}, x_{1}, \colr, t_{0})$. The excess is
deliberate and is the object of study rather than a free choice: the parameters are
introduced one at a time along a ladder, each rung scored on held-out Hubble residual
scatter under cross-validation and never on in-sample fit, with \salt\ fitted
independently on the same objects as an external benchmark. The ladder, which is also
the training schedule, is specified in the project's engineering record rather than
here.

The base rung is worth naming, because it is a controlled experiment rather than a
starting point. Rung L0 carries $(\dm, \colr, \tscale, \tmax)$ --- four parameters,
the same count as \salt, and in one-to-one correspondence with them. It differs in
\emph{mechanism} alone: where \salt\ produces stretch-like variation from an additive
component, L0 produces it from an exact reparameterisation. L0 against \salt\ is
therefore a like-for-like test of the mechanism at fixed parameter count, before any
new freedom is introduced. Everything above L0 asks whether additional freedom earns
its place.

\subsection{Gauge, assumption, prediction}
\label{sec:gauge}

Three kinds of condition appear above and must not be confused. A \emph{gauge}
condition removes a mathematical redundancy: predictions are identical before and
after, so it costs nothing, cannot be wrong, and is not testable. An \emph{identifying
assumption} is a statement about the world that pins down a split the data leave open;
it is substantive and may be false. A \emph{prediction} is neither, and is checked.

\begin{center}
\begin{tabular}{lll}
\toprule
Condition & Kind & Fixed by \\
\midrule
Arclength origin $\arc = 0$      & gauge      & the principal $g$-maximum, \eqref{eq:frameanchor} \\
Frame orientation                & gauge      & \eqref{eq:frameanchor}, for $n = 2$; a global \\
                                 &            & \emph{parameter} for $n \ge 3$ (two for $n=3$) \\
Metric $\mathrm{diag}(1,1)$      & gauge      & convention (Section~\ref{sec:magspace}) \\
Latent mean, covariance, sign    & gauge      & zero mean, unit diagonal covariance; $\shapeone$ \\
                                 &            & positively correlated with light-curve width \\
Position $\parallel (1,1)$       & gauge      & $\curve(0) = 0$, so $\dm$ is the peak $g$ magnitude \\
Position along $\hat{\mathbf{e}}_{c}$ & \emph{parameter} & $\colr$, free per supernova \eqref{eq:transcol} \\
Reddening direction $\hat{\mathbf{e}}_{c}$ & gauge & sample-wide unit vector; a \emph{parameter} \\
                                 &            & only at L2c; $(1,-1)/\sqrt{2}$ nested \\
Dust against shape               & \textbf{assumption} & \eqref{eq:idcond} \\
Colour \emph{evolution}          & prediction & terminal minus early offset $\perp (1,1)$; \\
                                 &            & the zero point is $\colr$, so only the difference \\
                                 &            & is predicted, not either colour absolutely \\
Total turning, terminal colour   & prediction & $\int\curv\,\dd\arc \approx \pi$, outside \eqref{eq:domain}; \\
                                 &            & qualitative, never imposed \\
\bottomrule
\end{tabular}
\end{center}

Two entries repay a second look, because both are places where a parameter was previously
counted where none exists. Orientation is gauge for $n = 2$: the rotational freedom and
the redundancy between the network's arclength origin and the frame cancel exactly, so an
orientation parameter would be a coordinate on that redundancy rather than added capacity
--- which is also why $n \ge 3$, where they do not cancel, genuinely carries
$n(n-1)/2 - 1$ of them. And $\hat{\mathbf{e}}_{c}$ is a gauge choice wherever $\colr$
multiplies a constant, since $\dm$ and $\colr$ then span $\Rtwo$ whatever direction is
chosen; only L2c's phase-varying $\delta\dustdir$ gives it content. Neither reading is
available from the parameter count alone; both follow from asking what the likelihood is
invariant under.

One substantive degeneracy of the timing layer remains.
A reparameterisation cannot change the curve in principle, but over a
finite window with noisy sampling a warp in \eqref{eq:reparam} and a change in $\curv$
can mimic one another. Regularise $\warpone, \warptwo$ toward zero and report the
fitted correlation between them and $\shape$. If it is large, the shape--timing
separation that motivates the model is not being realised, and that is a negative
result to report rather than absorb.

\section{Why the curve lives in magnitude space}
\label{sec:magspace}

The case is physical. Several of the quantities standing between an observed \snia\
and its intrinsic behaviour are, in magnitudes, simply \emph{added} to the
photometry --- and curvature does not respond to an added constant, since adding one
slides the curve bodily through the plane without bending it.

\begin{center}
\begin{tabular}{lll}
\toprule
Physical effect & Acts on the curve as & Removed? \\
\midrule
Luminosity distance        & rigid displacement $\parallel (1,1)$ & exactly \\
Absolute luminosity        & rigid displacement $\parallel (1,1)$ & exactly \\
Peculiar velocity          & rigid displacement $\parallel (1,1)$ & exactly \\
Per-band calibration error & rigid displacement, per band         & exactly \\
Time dilation, stretch     & reparameterisation, no displacement  & exactly \\
\midrule
Host and Galactic dust     & deformation: rigid part $+$ drift    & rigid part only \\
Intrinsic diversity \emph{(signal)} & deformation                 & no, and must not be \\
\bottomrule
\end{tabular}
\end{center}

The horizontal rule is the content of this section. Above it sit effects the geometry
disposes of for nothing; below it sit effects it cannot touch, one of which is the
signal.

\paragraph{The first variation of $\curv$ makes the rule precise.} Displace the curve
by a vector field decomposed in the Frenet frame,
$\curve \mapsto \curve + \varphi\,\tang + \psi\,\norml$. For a plane curve the induced
change in the signed curvature is, to first order,
\begin{equation}
  \Delta\curv = \psi'' + \curv^{2}\psi + \curv'\varphi ,
  \label{eq:firstvar}
\end{equation}
where $'$ denotes $\dd/\dd\arc$. Let the displacement be a \emph{constant} vector
$\mathbf{v}$, so $\varphi = \tang\cdot\mathbf{v}$ and $\psi = \norml\cdot\mathbf{v}$.
Then $\psi' = -\curv\varphi$ and $\varphi' = \curv\psi$ by \eqref{eq:frenet2d}, so
$\psi'' = -\curv'\varphi - \curv^{2}\psi$ and
\begin{equation}
  \Delta\curv = \bigl(-\curv'\varphi - \curv^{2}\psi\bigr) + \curv^{2}\psi
                + \curv'\varphi = 0
  \label{eq:transinv}
\end{equation}
identically, for every curve and every constant $\mathbf{v}$. That is the row group
above the rule, and its whole content: what matters is that $\mathbf{v}$ be
independent of $\arc$, not that it be achromatic. A zeropoint error differs in every
band, is strongly chromatic, and is still annihilated.

Conversely nothing in \eqref{eq:firstvar} vanishes when the displacement depends on
$\arc$. Dust is such a displacement, because the absorption a broad filter delivers
depends on the spectrum inside it and that spectrum evolves
(Section~\ref{sec:phasedust}); so is intrinsic diversity, which is the signal
$\curv(\arc;\shape)$ exists to describe. Dust and signal enter $\curv$ through the
same channel, and no choice of ambient space or invariant separates them --- separation
can come only from population-level assumptions, stated as such in
Sections~\ref{sec:phasedust} and~\ref{sec:gauge}. What magnitude space buys is
narrower and still decisive: by \eqref{eq:transinv} the quantity being inferred is
\emph{exactly} orthogonal to shape, which is why $\dm$ may be free per supernova and
absorb peculiar velocity exactly. Being the same type as the signal is not being the
same size; Section~\ref{sec:phasedust} bounds the size using \eqref{eq:firstvar}
itself.

Equation~\eqref{eq:transinv} also bounds what the geometry delivers: $\curv$ is
insensitive to every constant offset, so it carries no information about absolute
magnitude or absolute colour, only about their evolution. Those are supplied
separately, by $\dm$, by the boundary conditions of Section~\ref{sec:place}, and by
$\colr$ in rung L2c.

\paragraph{In flux the same effects are multiplicative.} Distance is an isotropic
dilation $f \mapsto af$, giving $\arc \mapsto a\arc$ and $\curv \mapsto \curv/a$: a
photographic enlargement, so shape survives but $\curv$ does not, and the damage is
repairable in principle by a scale-invariant combination such as
$\arclen\,\curv(\arc/\arclen)$. Dust is an \emph{anisotropic} dilation,
\begin{equation}
  (\fluxg, \fluxr) \mapsto
  \bigl(\fluxg\,10^{-0.4 A_{g}},\ \fluxr\,10^{-0.4 A_{r}}\bigr),
  \qquad A_{g} \neq A_{r},
\end{equation}
which takes a circle to an ellipse and constant curvature to non-constant; no
rescaling or reparameterisation recovers it. Dust, not distance, settles the choice of
space. The likelihood is nevertheless evaluated in native flux
(Section~\ref{sec:likelihood}), for two reasons independent of geometry: \ztf\
difference-imaging errors are approximately Gaussian in flux and negative fluxes are
meaningful. \textbf{No data are ever converted to magnitudes; only the
model is expressed there.}

\paragraph{No per-supernova effect acts as a rotation.}
\label{sec:rotation}
Every entry in the table above is a
translation, a phase-dependent displacement, or a reparameterisation; none is a
rotation. A per-supernova rotation would be a degree of freedom with no physical
referent, and none is fitted at any rung. A supernova-dependent $R_{V}$ is not a
counterexample: it changes the \emph{direction} of the reddening displacement,
selecting a different $\dustdir(\arc)$ in \eqref{eq:dustcurve}, but does not rotate the
curve. The template as a whole is the exception, and only as a whole: it carries one
orientation that no physical effect determines --- the rigid freedom left by
\eqref{eq:transinv} --- and for $n = 2$ even that is fixed by convention rather than
fitted, since \eqref{eq:frameanchor} exhausts it; for $n \ge 3$ what remains is a single
rotation shared across the sample. Neither is a per-supernova freedom. The reddening
direction $\hat{\mathbf{e}}_{c}$ is sample-wide for the same reason: a per-supernova
reddening direction would rotate the colour axis differently for each object with no
physical referent.

\paragraph{The metric is a convention.} Curvature has dimensions of inverse length and
so presupposes a metric on magnitude space, of which none is physically preferred:
nothing dictates how one magnitude of $g$ compares in length to one magnitude of $r$.
Calling $(1,1)/\sqrt{2}$ and $(1,-1)/\sqrt{2}$ orthogonal unit vectors presumes
$\mathrm{diag}(1,1)$. The arguments of this section are untouched by the choice, since
translations are isometries of every translation-invariant metric; the numerical values
of $\curv$ are not, since rescaling $\magr \mapsto 2\magr$ alters every curvature while
altering nothing physical. Two further facts sharpen \eqref{eq:transinv}: among maps of
the ambient space, isometries are the \emph{only} curvature-preserving ones, which is
what makes the flux comparison a derivation rather than a preference; and the signed
curvature of \eqref{eq:frenet2d} picks up $\det Q$ under an isometry of linear part
$Q$, so the invariance group is $\SE(n)$ and only $\abs{\curv}$ is
reflection-invariant. Normalisation convention: $\dm$ is the offset added to
\emph{each} band, as in \eqref{eq:forward}, so the displacements
\begin{align}
  \text{distance, luminosity, peculiar velocity:} \quad
    &\curve \mapsto \curve + \dm\,(1,1), \label{eq:transdist}\\
  \text{reddening (phase-independent part):} \quad
    &\curve \mapsto \curve + \colr\,\hat{\mathbf{e}}_{c}, \label{eq:transcol}
\end{align}
with $\hat{\mathbf{e}}_{c}$ the fitted global reddening direction of
Section~\ref{sec:place}, a unit vector whose fixed-$R_{V}$ idealisation is
$(1,-1)/\sqrt{2}$. The distance displacement has length $\dm\sqrt{2}$, since $(1,1)$ is
not normalised; a stray $\sqrt{2}$ from that convention would propagate into the
absolute-magnitude offset of \eqref{eq:tripp}.

\section[Dust extinction is phase-dependent]{Dust extinction is phase-dependent, so
\eqref{eq:transcol} is only approximate}
\label{sec:phasedust}

Equation~\eqref{eq:transcol} treats reddening as a fixed vector. It is not.
Broadband extinction is a flux-weighted average of the monochromatic law over the
filter,
\begin{equation}
  A_{X}(\phase) = -2.5\log_{10}
    \frac{\int S_{X}(\lambda)\,F(\lambda,\phase)\,10^{-0.4A(\lambda)}\,\dd\lambda}
         {\int S_{X}(\lambda)\,F(\lambda,\phase)\,\dd\lambda},
  \label{eq:bandext}
\end{equation}
and the weights are the supernova's own spectral energy distribution
$F(\lambda,\phase)$, which evolves strongly over the event \citep{Nugent2002}. As the
\snia\ reddens intrinsically the flux-weighted effective wavelength within each filter
shifts, so the extinction suffered per unit reddening drifts with phase: both the
magnitude and the direction of the displacement change. This is why
\eqref{eq:dustcurve} is written with an $\arc$-dependent $\dustdir$: dust acts not as a
translation but as a \emph{phase-dependent displacement}, which is not an isometry and
does change both $\curv$ and arclength.

\paragraph{Order counting: the shape contamination is second order.} Split the
displacement into its phase-average and the variation about it,
\begin{equation}
  \colr\,\dustdir(\arc)
    = \underbrace{\colr\,\bar{\dustdir}}_{\text{translation}}
    + \underbrace{\colr\,\delta\dustdir(\arc)}_{\text{deformation}},
  \qquad
  \bar{\dustdir} = \langle \dustdir \rangle,
  \quad \langle \delta\dustdir \rangle = 0 .
  \label{eq:dustsplit}
\end{equation}
Substituting into \eqref{eq:firstvar} and using \eqref{eq:transinv} on the constant
part, the induced curvature change is
\begin{equation}
  \Delta\curv(\arc) = \colr\Bigl[
      \bigl(\norml\cdot\delta\dustdir\bigr)''
    + \curv^{2}\bigl(\norml\cdot\delta\dustdir\bigr)
    + \curv'\bigl(\tang\cdot\delta\dustdir\bigr)\Bigr] ,
  \label{eq:dkappa}
\end{equation}
in which $\bar{\dustdir}$ does not appear at all. Writing
$\varepsilon = \norm{\delta\dustdir}/\norm{\bar{\dustdir}}$ for the fractional phase
variation and $\ell$ for the arclength scale over which it varies,
\begin{equation}
  \abs{\Delta\curv} = \mathcal{O}\!\left(
    \colr\,\varepsilon\,\norm{\bar{\dustdir}}\,
    \max\!\bigl(\ell^{-2},\,\curv^{2}\bigr)\right),
  \label{eq:orders}
\end{equation}
so the contamination of shape is small twice over: once in the reddening amplitude
$\colr$, and once in $\varepsilon$, which is small because $A_{X}$ in
\eqref{eq:bandext} responds only weakly to SED evolution inside a broad filter --- the
flux-weighted effective wavelength shifts by a small fraction of the filter width.
Intrinsic diversity, by contrast, enters $\curv$ at first order, being what
$\curv(\arc;\shape)$ parameterises. The expectation is thus a \textbf{hierarchy} ---
the phase-dependent variance the geometry sees should be dominated by intrinsic
supernova dispersion --- and it is an expectation to be measured, not assumed.

Two limits on the argument. It is amplitude-dependent through the $\colr$ in
\eqref{eq:dkappa}, so it degrades for the reddest objects and every diagnostic below is
reported \emph{stratified by} fitted $\colr$ rather than sample-averaged. And it covers
the \emph{shape} channel only: in the colour channel $\colr\,\bar{\dustdir}$ and a
per-supernova intrinsic colour offset are the same displacement, so they mix at first
order, exactly as \salt's $\colr$ does. That degeneracy is not resolved here.

This is why L2c fits the amplitude of $\dustdir(\arc)$ rather than of a constant
vector: by \eqref{eq:dustsplit} a constant would remove $\colr\,\bar{\dustdir}$ and
leave $\colr\,\delta\dustdir$ free to be absorbed by $\shape$, which is the
contamination the rung exists to prevent, and the parameter count is the same either
way. Below that rung the only
ambient motion is $\dm$, so the template must predict each supernova's absolute colour
from its shape alone and host dust, having nothing else to occupy, is forced into
$\shapeone, \shapetwo$; whether shape can absorb reddening or needs its own direction
is a result either way.

\paragraph{Dust against shape, and the identifying assumption.} $\colr$ is degenerate
against $\shape$ to the extent that $\delta\dustdir$ lies in
$\mathrm{span}\{\partial\curve/\partial\shape\}$ at the observed epochs. This is not a
redundancy that can be gauged away but a genuine ambiguity of physical cause. The
identifying assumption is
\begin{equation}
  \mathrm{Corr}(\colr, \shape) = 0
  \quad\text{over the population,}
  \label{eq:idcond}
\end{equation}
justified by the physical independence of line-of-sight dust from explosion physics and
known to be only approximate --- the host mass step couples \snia\ properties to host
environment, and dust column is environmental. By \eqref{eq:orders} it is expected to
function as a \emph{check} rather than to carry weight, so it is imposed \emph{and} its
unimposed value reported. It constrains only the interpretation of the split, so L2c's
held-out Hubble score is meaningful either way.

\paragraph{Diagnostics.} Three numbers, in increasing directness.

\begin{enumerate}
  \item $\mathrm{Corr}(\colr, \shape)$ over the population, computed both with and
        without \eqref{eq:idcond} imposed. A large unconstrained value falsifies
        \eqref{eq:orders}.
  \item The error-weighted overlap of the dust direction with the shape directions at
        the epochs actually observed. Let $\mathbf{d}$ stack
        $\partial f_{X}/\partial\colr$ over observed epochs, $J$ have columns
        $\partial f_{X}/\partial\shapeone$, $\partial f_{X}/\partial\shapetwo$,
        $W = \mathrm{diag}(1/\sigma_{i}^{2})$, and
        $P = J(J^{\!\top} W J)^{-1}J^{\!\top}W$. Then
        \begin{equation}
          \rho^{2} = \frac{\mathbf{d}^{\!\top} W P\, \mathbf{d}}
                          {\mathbf{d}^{\!\top} W \mathbf{d}} \in [0,1]
          \label{eq:overlap}
        \end{equation}
        is the fraction of the dust signal this cadence and these error bars cannot
        distinguish from shape. Two details decide whether it is meaningful. It is
        built in \emph{flux}, where the errors are Gaussian, so
        $\partial f_{X}/\partial\shape = -0.4\ln\!10\,f_{X}\,
        \partial m_{X}/\partial\shape$; and $\mathbf{d}$ is formed from
        $\delta\dustdir$, not $\dustdir$, since by \eqref{eq:transinv} the mean part
        cannot be confused with shape and including it would inflate $\rho^{2}$.
  \item $\Delta\curv$ from \eqref{eq:dkappa}, with $\dustdir$ built by evaluating
        \eqref{eq:bandext} on a grid of $E(B-V)$ and $R_{V}$, against the range of
        $\curv$ the network spans over fitted $\shape$. This tests \eqref{eq:orders}
        directly. Small at every $\colr$ licenses the interpretation of $\shape$; order
        unity at any $\colr$ means the shape latents are substantially measuring dust
        there --- a negative result for interpretation, to be reported as one, and not
        an obstacle to the Hubble-diagram result, for which the causes need not be
        separated.
\end{enumerate}

The external SED template \citep{Hsiao2007} is used here only to quantify a systematic,
never to fit, initialise, or K-correct, so independence from \salt\ is preserved. The discussion
behind several points compressed above --- the tolerability of the mixing for the
distance indicator, the status of the identifying assumption, the cost of holding
$\dustdir$ fixed, and why $\zred < 0.05$ does not help with dust --- is held in
\texttt{docs/model-note-backlog.md} until this note stabilises.

\section{Samples}
\label{sec:samples}

Coverage was measured directly from the \ztf\ SN~Ia DR2 archive
\citep{Rigault2025}. Starting from the 3628 catalogued objects, we require
$\zred < 0.05$ and both the \texttt{lccoverage} and \texttt{fitquality} flags, leaving
669 supernovae, of which 654 have a light-curve file present in the release. Epochs are
counted inside the rest-frame window $[-15,+40]$~d with signal-to-noise above five and
with none of the quality bits $\{1,2,4,8,16\}$ set.

\begin{center}
\begin{tabular}{lrrrrr}
\toprule
Minimum epochs per band & $g$ & $r$ & $i$ & $g\,\&\,r$ & $g\,\&\,r\,\&\,i$ \\
\midrule
$\geq 3$  & 636 & 639 & 216 & 636 & 215 \\
$\geq 5$  & 606 & 631 & 179 & 599 & 177 \\
$\geq 10$ & 490 & 537 &  82 & 465 &  80 \\
\bottomrule
\end{tabular}
\end{center}

Two samples follow. The \textbf{primary sample} is the 599 supernovae with at least
five good $g$ and $r$ epochs; here $n=2$ and there is a single invariant. The
\textbf{torsion subsample} is the 177 that additionally have at least five good
$i$-band epochs; here $n=3$ and $\tors(\arc)$ exists as a genuine second invariant.
The latter is small, but it is the only configuration in which the original motivation
for a second shape parameter --- that torsion describes the evolution of inter-band
magnitude differences --- can be tested, so it is a distinct analysis rather than an
afterthought. Torsion is undefined wherever $\curv_{1}$ vanishes
(Section~\ref{sec:curve}), so it is identified only where the track turns.

\paragraph{A caveat on the window used for counting.} Counting epochs inside
$[-15,+40]$~d requires a date of maximum, but $\tmax$ is precisely what the model fits
(Section~\ref{sec:traversal}); the numbers above were obtained using the archive's
\salt\ $t_{0}$ to place the window. This is \salt\ entering upstream of
\emph{selection}, far weaker than \salt\ entering the fit --- no photometry is
transformed, initialised, or K-corrected by it, and the fitted $\tmax$ remains free and
independent. It is nonetheless an exception to the rule that \salt\ stays strictly
downstream, and is recorded here rather than left implicit. The check that settles it is
to recount using an SED-free anchor, for instance a window placed on the brightest good
$g$ epoch, and confirm the sample sizes move negligibly. The 36 objects with NaN
$t_{0}$ must in any case be selected by such an estimator or excluded explicitly.

\paragraph{Milky Way extinction.} Galactic reddening is tabulated per object as
\texttt{mwebv}, so its \emph{amplitude} is known and is removed deterministically
before fitting rather than being fitted. Only host-galaxy reddening remains a free
quantity, and only in rung L2c. Note that what is known is $E(B-V)$, not
$A_{g}(\phase)$ and $A_{r}(\phase)$; by Section~\ref{sec:phasedust} the conversion is
phase-dependent and needs an SED, so removing a constant offset is an approximation
here as well --- milder, since Galactic columns at the high latitudes of this sample
are low, but the same approximation and not exact.

\section{Relation to \salt, and distance}
\label{sec:salt}

The two models can be placed side by side once the geometric quantities are mapped
onto their photometric counterparts.

\begin{center}
\begin{tabular}{lll}
\toprule
Concept & \salt & This work \\
\midrule
Underlying object & SED surface $F(p,\lambda)$ & curve $\curve(\arc) \subset \Rn$ \\
Basis             & B-splines in $(p,\lambda)$ & neural network in $\arc$ \\
Stretch           & \emph{additive} $x_{1}M_{1}(p,\lambda)$ & \emph{reparameterisation} $\tscale$ \\
Residual shape    & --- (same $x_{1}$)          & $\shapeone$ conditioning $\curv$ \\
New freedom       & ---                         & $\shapetwo$ conditioning $\curv$ \\
Timing origin     & $t_{0}$                     & $\tmax$; warp $\warpone, \warptwo$ \\
Normalization     & $x_{0}$                     & $\dm$ \\
Colour            & $\colr$ with law $\mathrm{CL}(\lambda)$ & $\colr$, translation; $\colr\,\dustdir(\arc)$ at L2c \\
K-corrections     & built in via the SED        & suppressed by $\zred < 0.05$, not removed \\
\bottomrule
\end{tabular}
\end{center}

The decisive structural difference is the treatment of timing, and it is sharper than
``\salt\ has one time parameter''. \salt\ has \emph{no} mechanism for
reparameterising time at all. A stretch $t \mapsto t/\stretchsym$ leaves the curve
$\curve$ pointwise unchanged and alters only its traversal; \salt\ instead represents
stretch-like variation by the \emph{additive} component $x_{1}M_{1}(p,\lambda)$, which
is a different operation. The two agree only to first order,
\begin{equation}
  m\!\left(\phase/\stretchsym\right)
    = m(\phase) - (\stretchsym - 1)\,\phase\, m'(\phase)
      + \mathcal{O}\!\left((\stretchsym-1)^{2}\right),
  \label{eq:stretchlin}
\end{equation}
so a linear component reproduces a stretch only in the small-stretch limit, with
$M_{1} \propto -\phase\,\partial M_{0}/\partial \phase$. This is not a reconstruction
after the fact: \salt\ initialises $M_{1}$ as exactly such a finite difference, the
SED sequence at $\stretchsym = 1.1$ minus that at $\stretchsym = 1$
\citep{Guy2007}. The component is then fitted freely and does not remain a pure
linearised stretch, but the residual nonlinearity is visible in the published
$x_{1} \to \stretchsym$ conversion, which requires a cubic.

Two consequences. First, \salt\ reports that most of the variability $x_{1}$ captures
\emph{is} a stretch, so the dominant mode of \snia\ diversity is precisely the one
\eqref{eq:reparam} represents exactly and \eqref{eq:stretchlin} shows \salt\
linearises. That is the strongest structural argument for the present model. Second,
the same fact predicts that $\shapeone$ should be \emph{small}: with stretch removed
into $\tscale$, what remains for the first latent is only shape variation no
reparameterisation can produce. A concrete check follows --- fitted $\tscale$ should
track the published conversion
$\stretchsym = 0.98 + 0.091x_{1} + 0.003x_{1}^{2} - 0.00075x_{1}^{3}$, and the scatter
about it measures what the linearisation misses.

On colour the comparison runs the other way, and the point is worth conceding plainly.
\salt's colour law $\mathrm{CL}(\lambda)$ is phase-independent, but \salt\ carries an
SED surface from which \eqref{eq:bandext} is evaluated at each phase, so the phase
dependence emerges automatically; having no SED, this model must import
$\dustdir(\arc)$ from an external template. \salt\ is structurally better placed here.
What the two share is the mixing itself: \salt's $\colr$ conflates dust with intrinsic
colour, which is why $\beta$ is fitted rather than predicted, the direct analogue of the
global reddening direction $\hat{\mathbf{e}}_{c}$ fitted here.

Two further costs. Having no SED, the redshift cut suppresses only the K-correction
branch of the missing-SED problem \citep{Nugent2002}: the phase dependence of
\eqref{eq:bandext} is a rest-frame effect, identical at $\zred = 0$, so
Section~\ref{sec:phasedust} is a mitigation and not a solution. And full conditioning
in \eqref{eq:kappanet} means the latents are not linear coefficients, so ablation is
less clean than in \salt\ and the gauge fixing of Section~\ref{sec:gauge} is a partial
remedy only.

Distance is not delivered directly by the geometry. Since $\dm$ is free per supernova,
the geometry supplies standardization \emph{variables}, and distance follows from a
Tripp-style relation \citep{Tripp1998},
\begin{equation}
  \dm_{\mathrm{corr}} = \dm - \alpha\,\shapeone - \beta\,\colr
                        - \gamma\,\shapetwo + \mathcal{M},
  \label{eq:tripp}
\end{equation}
with $\alpha$, $\beta$, $\gamma$ and $\mathcal{M}$ fitted globally. The first three
terms are \salt's own standardisation in one-to-one correspondence --- $\shapeone$ for
$x_{1}$, $\colr$ for $\colr$ --- so the comparison is direct. The fourth is the whole
point: $\gamma$ is the coefficient of the one degree of freedom \salt\ does not have,
and whether it is nonzero, and whether it reduces held-out scatter, is the headline
result the project is built to produce. It is also why the
dust--shape mixing of Section~\ref{sec:phasedust} is asymmetric in its consequences:
for the distance indicator the mixture is absorbed by the fitted $\alpha, \beta$ and
costs nothing, whereas for reading $\shape$ as explosion physics it is a limitation
whose size must be measured and stratified by $\colr$. The two conclusions are reported
separately.

\bibliographystyle{plainnat}
\bibliography{refs}

\end{document}
